\documentclass[conference]{IEEEtran}
\usepackage[utf8]{inputenc}
\usepackage[T1]{fontenc}
\usepackage{cite}
\usepackage{amsmath,amssymb}
\usepackage{graphicx}
\usepackage{booktabs}
\usepackage{url}

\title{Effectiveness of a Generate--Verify--Repair Pipeline with Batfish for LLM‑Generated Vendor CLI Configurations}

\author{
\IEEEauthorblockN{Autonomous AI Researcher}
\IEEEauthorblockA{CNSM 2026 Agentic AI Researcher Track}
}

\begin{document}

\maketitle

\begin{abstract}
We report a preregistered paired experiment (cnsms2026-001-gvr-batfish-prereg-v1) that evaluated whether a generate--verify--repair (GVR) pipeline---shared initial LLM candidate (gpt-5-mini) + a deterministic Batfish‑style validator + at most one automated LLM repair---reduces deterministic‑verifier detected semantic/safety failures on intent\ensuremath{\rightarrow}vendor‑CLI tasks. Planned N=300 pairs; 266 complete pairs were analysed after preregistered exclusions. Baseline pass-rate = 260/266 (97.744\%); guarded (final GVR) pass-rate = 265/266 (99.624\%). Paired difference = 0.01879699248; exact McNemar two‑sided p = 0.0625. Paired bootstrap (B=10,000, seed=1729) 95\% percentile CI = [0.0037593984962406013, 0.03759398496240601]. Repair was invoked for 6 tasks; median repair iterations per task = 0. Per‑episode latency was not recorded in per‑task artifacts and thus could not be evaluated. Under shared\_initial\_candidate semantics the observed absolute improvement is small and did not meet the preregistered \ensuremath{\alpha}=0.05 decision rule. We provide preregistered design details, deterministic validator semantics, exact paired counts, representative artifact‑grounded repair examples, contamination findings (no exact public duplicates flagged), and a focused discussion of limitations and implications.
\end{abstract}

\section{Introduction}
Generative large language models (LLMs) are actively explored for NetOps tasks such as intent\ensuremath{\rightarrow}CLI translation, zero‑touch configuration synthesis, and automated maintenance workflows [1], and surveys emphasize both opportunity and risk when LLMs are applied to operational network configuration [2]. Stochastic LLM outputs can produce syntactic or semantic violations---missing required commands, unsafe command orderings, or constraint breaches---that create operational risk if applied directly. To mitigate that risk, pipelines that interpose deterministic verifiers (e.g., Batfish or header‑space analyses) and bounded repair loops---collectively generate--verify--repair (GVR)---have been proposed to provide deterministic acceptance criteria and automated correction [3,4]. Prior work presents components and prototypes, but preregistered, paired evaluations that archive verifier traces and quantify verifier‑triggered repair benefit on reproducible NetOps task banks remain limited.

This study tested the preregistered question (cnsms2026-001-gvr-batfish-prereg-v1): Does a GVR pipeline that uses a single sampled initial LLM candidate per task (shared\_initial\_candidate semantics), a deterministic Batfish‑style validator, and at most one automated LLM repair call reduce deterministic‑verifier detected semantic/safety failures relative to the same initial candidate without repair? Our principal contribution is a preregistered, artifact‑archived paired evaluation executed under the CNSM Agentic AI Researcher constraints that reports exact paired counts, inferential outcomes, execution accounting, representative artifact examples, and an evidence‑grounded interpretation.

\section{Related Work}
Prior work relevant to this study spans (a) applied LLM prototypes and benchmarks for network configuration, (b) surveys documenting LLM failure modes and recommendations to ground generation, (c) proposals for generate--verify--repair and related architectures for deterministic validation, and (d) established deterministic network verification tools that serve as practical oracles in automated pipelines. We synthesize those threads and emphasize aspects most relevant to a preregistered paired GVR evaluation.

LLM-driven NetOps prototypes and benchmarks: Recent system‑level and benchmark efforts investigate whether LLMs can aid network configuration generation and automation. NetConfEval and similar task collections examine intent\ensuremath{\rightarrow}CLI translation and measure correctness under vendor‑dialect constraints; these works provide public task exemplars and motivate reproducible task generators for controlled evaluation [1]. Such benchmark tasks typically encode required command sequences and vendor syntax, which directly map to the validator check types we archive in this study (required‑sequence, allowed‑field checks, topology reachability). Importantly, prior benchmarks often evaluate single‑pass generation accuracy rather than integrated verifier‑triggered repair, which motivates the paired GVR design adopted here.

Failure modes, grounding, and mitigation recommendations: Survey and synthesis papers summarize that LLMs commonly produce both syntactic and semantic errors when applied to domain‑constrained tasks: missing commands, incorrect ordering, and hallucinated or unsafe modifications [2]. Those surveys recommend grounding strategies (retrieval, tool use), constrained generation, and verifier insertion to reduce the operational risk of accepting model outputs without deterministic checks. Our work adopts those prescriptions by interposing a deterministic validator and bounding the repair loop to one automated correction attempt, aligning experimentally with the practical mitigations advocated in the literature.

Generate--verify--repair and bounded repair proposals: The idea of interposing deterministic oracles and using closed‑loop repair is well established in proposals for regulated or safety‑sensitive domains. Convergent GVR architectures aim to combine stochastic generation with deterministic acceptance criteria and limited repair to achieve practical correctness guarantees while preserving the generative model's productivity [3]. These proposals emphasize publishing verifier traces and repair prompts so that corrected outputs are auditable---an archival design we implemented in the execution artifacts (validator traces, repair\_prompt\_sha256, repair provider traces).

Deterministic network verifiers as oracles and their coverage limits: Tools like Batfish and header‑space analyses provide deterministic, rule‑based checks for reachability, ordering, and configuration invariants and have been incorporated into research prototypes that integrate model outputs with formal validation [4]. Batfish‑style validators excel at expressing and mechanically checking reachability and required‑sequence constraints on modeled topologies; however, their guarantees are strictly limited to the rule battery they encode. This distinction matters experimentally: a verifier that checks an explicit finite set of syntactic/semantic invariants yields a binary pass/fail endpoint that is reproducible and auditable, but it cannot detect semantic failures outside its coverage. Our deterministic\_netops\_validator\_v5 follows this same engineering trade‑off (syntactic parsing, required‑sequence checks, allowed‑fields enforcement, and topology reachability), and we archive per‑episode validation\_before/after traces so readers can inspect exactly which rules were exercised on each episode.

What LLMs need to synthesize correct router configurations and constrained decoding: Empirical studies analyzing LLM difficulty with router configuration point to sequencing, vendor‑specific syntax, and dependency constraints as recurring challenges [5]. These analyses motivate two complementary mitigations: (1) stronger first‑pass constraints or constrained decoding to reduce syntactic/format errors, and (2) verifier‑triggered repair to correct semantic violations that slip through first pass. Our study deliberately fixes the initial candidate (shared\_initial\_candidate) and thereby isolates the marginal effect of the verifier‑triggered repair step rather than conflating improvements due to constrained first‑pass prompting or multiple draws.

Gaps in prior evaluations that motivated this study: The supplied literature shows active proposals and component evaluations but comparatively fewer preregistered, paired experiments that (i) use shared\_initial\_candidate semantics to measure the isolated effect of bounded repair, (ii) archive raw model traces together with deterministic validator traces and repair provider traces for each episode, and (iii) apply explicit exclusion, retry, and contamination rules in a confirmatory design. Those gaps motivated our preregistered paired design and the artifact‑heavy archival choices (manifested in the execution and analysis artifacts we publish alongside this manuscript).

Methodologically relevant syntheses for experimental design: From the surveyed and prototyped prior work we derived three concrete design choices: (1) use a deterministic validator battery (Batfish‑style) to obtain an auditable binary endpoint; (2) bound the repair loop (<=1 automated repair call per failed candidate) to maintain a clear paired estimand under shared\_initial\_candidate semantics; and (3) preregister explicit missingness, retry, and contamination treatments to avoid post hoc analytic choices. Each of these choices has precedent or justification in the cited literature: verifiers as deterministic oracles [4], GVR architectures for regulated settings [3], and benchmark/task construction recommendations that emphasize reproducibility and contamination controls [1,2].

In sum, our related‑work synthesis locates this study at the intersection of applied intent\ensuremath{\rightarrow}CLI benchmarks, survey‑driven mitigation recommendations, GVR methodological proposals, and deterministic verification tooling. The specific empirical contribution is a preregistered, archived paired GVR evaluation under shared\_initial\_candidate semantics that fills a documented gap in artifact‑grounded confirmatory assessments of verifier‑triggered repair effectiveness.

\section{Methodology}
Preregistration and estimand: The experiment followed the preregistered plan cnsms2026-001-gvr-batfish-prereg-v1 (manifest SHA‑256: 46a22349\allowbreak{}2f760950\allowbreak{}b3d06f47\allowbreak{}5129ce21\allowbreak{}e52e296c\allowbreak{}4c2a1e09\allowbreak{}df3ccd50\allowbreak{}bec9f891). The primary estimand is the paired\_success\_rate\_difference\_guarded\_minus\_baseline under shared\_initial\_candidate semantics. Under these semantics the identical single sampled initial model candidate per task is used to produce both outcomes: (1) baseline --- score the single initial candidate with the deterministic validator; (2) guarded --- if that same candidate fails, the pipeline may invoke at most one automated LLM repair call and then rescore with the same deterministic validator. This design isolates the verifier‑triggered repair effect for a fixed initial draw and matches the preregistered execution\_contract (generation\_semantics = "shared\_initial\_candidate").

Task bank, sampling, and inclusion rules: The task bank was deterministically generated to produce N=300 intent\ensuremath{\rightarrow}CLI pairs, stratified into two vendor‑dialect clusters (150 Cisco‑like, 150 Juniper‑like) via a seeded synthetic generator supplemented with public NetConfEval‑style items as specified in the preregistration. Generator ids and deterministic seeds are archived in the task manifest. Inclusion and exclusion rules followed the preregistered contract: public exact‑duplicate tasks detected by the contamination check were to be excluded from the primary confirmatory analysis (none were excluded for exact duplication in this run); provider/API transient failures triggered up to two automatic retries, and pairs where both arms failed due to infrastructure/model API errors were excluded from the primary paired analysis per the 'complete\_pair\_primary' missingness rule. The execution manifest records planned\_episode\_count = 600 (300 pairs \ensuremath{\times} 2 arms) and completed\_episode\_count and exclusions are reconciled in the deterministic reconciliation artifacts.

Model orchestration and runtime configuration: The declared model used for generation and any repair calls was recorded as gpt-5-mini (execution/model\_configuration.json declared\_model\_version = "gpt-5-mini"). Archived configuration fields include max\_output\_tokens = 1024, maximum\_attempts\_per\_call = 1, provider = openai\_responses, and temperature = null (unset). Provider accounting recorded hosted\_model\_call\_event\_count = 306, completed\_hosted\_model\_call\_event\_count = 272, and failed\_hosted\_model\_call\_event\_count = 34; stage accounting recorded shared\_initial\_generation = 300 and repair = 6. Because temperature was unset and no centralized deterministic sampling seed for provider calls is archived, nondeterministic sampling of provider responses cannot be excluded for the recorded calls; this runtime nondeterminism is declared and discussed in Limitations and the Disclosure Statement.

Deterministic validator design and scoring rules: The binary oracle was deterministic\_netops\_validator\_v5. For each candidate it executed a fixed rule battery that is archived per episode: (1) CLI syntactic parsing and normalization; (2) required‑sequence checks comparing the task's required\_sequence to the parsed commands; (3) constrained‑field touch checks against allowed\_touched\_fields and allowed\_fields in the task payload; and (4) reachability/constraint validation over the deterministic synthetic topology model (a lightweight Batfish‑style evaluation). Each validator trace archives validation\_before and validation\_after objects containing parsed\_command\_count, required\_command\_count, observed\_touched\_field\_count, violation codes, normalized\_configuration, and a boolean valid flag. The confirmatory scoring rule for the primary binary endpoint was 'full\_intent\_constraint\_validation' as recorded in the scoring artifacts; these artifacts drive the paired contingency counts used in the primary test.

Repair loop mechanics and bounding: Per the preregistration the automated repair stage was strictly bounded to at most one LLM repair call per episode when the initial candidate failed the validator. The repair prompt construction is archived: it includes the task constraints, the validator failure codes and normalized validator diagnostics, and an explicit instruction to produce a corrected configuration satisfying the same task constraints. Repair provider traces and repair\_prompt identifiers are archived for each repair invocation. Stage accounting indicates 6 repair calls executed in the run; all repair artifacts (repair prompts, repaired outputs, and subsequent validator traces) are included in the evidence bundle and illustrated in representative scoring JSONs.

Analysis plan and inferential procedures: The prespecified primary hypothesis test was the exact McNemar two‑sided test on paired binary outcomes at \ensuremath{\alpha}=0.05 (estimator id paired\_success\_rate\_difference\_guarded\_minus\_baseline). Interval estimation used a paired bootstrap percentile CI (B=10,000, bootstrap\_seed = 1729) for the paired difference to quantify effect magnitude and uncertainty; these bootstrap parameters and resample counts are archived in the analysis log. Missingness handling for the primary analysis followed the preregistered 'complete\_pair\_primary' rule: pairs with both‑arm provider/infrastructure failures were excluded from the confirmatory McNemar test; a preregistered sensitivity analysis treats excluded episodes conservatively as failures. Secondary analyses (preregistered) included median\_repair\_iterations\_per\_task and median end‑to‑end latency; the latter could not be computed because per‑episode timing was not archived. Multiplicity control for formal secondary tests was specified with a Bonferroni adjustment; exploratory analyses were explicitly labeled and not used to support the primary claim.

Execution logging, archival, and auditability: Every episode archived raw model outputs, normalized configurations, validator\_before/after traces, scorer decision codes, and provider call accounting. Deterministic reconciliation artifacts verify episode accounting (completed\_episode\_count, failed\_episode\_count, complete\_pair\_count) and provider call tallies. Key execution statistics recorded in the reconciliation artifacts are: planned\_episode\_count = 600, completed\_episode\_count = 532, failed\_episode\_count = 68, complete\_pair\_count = 266, model\_calls\_used = 306, repair\_calls = 6, cache\_hit\_count = 0, cache\_miss\_count = 272. These archived artifacts permit independent verification of the paired counts (n11, n10, n01, n00) and replay of repaired episodes for failure‑mode inspection.

Design constraints and estimand implications: We emphasize that shared\_initial\_candidate semantics (one sampled initial candidate per task used to evaluate both baseline and guarded outcomes) was selected and executed exactly as preregistered. This choice isolates the incremental effect of a validator‑triggered, bounded repair for the observed initial sample; it does not measure the effect of constrained prompting, ensemble first‑pass generation, or multiple independent initial draws. Those alternative interventions would require independent-generation arms or additional preregistered conditions that were not executed under this run and therefore are outside the estimand and scope of the reported confirmatory analysis.

\section{Results}
Execution accounting and sample flow: The preregistered plan targeted 300 independent intent\ensuremath{\rightarrow}CLI tasks (600 episodes across two conditions). The run recorded planned\_episode\_count = 600 and completed\_episode\_count = 532; provider/infrastructure transient failures produced failed\_episode\_count = 68. Applying the preregistered 'complete\_pair\_primary' exclusion rules (exclude both‑arm provider/infrastructure failures) yielded complete paired units analysed = 266. The analysis pipeline's provider accounting indicates hosted\_model\_call\_event\_count = 306 total model‑call events across the run, with completed\_hosted\_model\_call\_event\_count = 272 and failed\_hosted\_model\_call\_event\_count = 34; stage accounting records shared\_initial\_generation = 300 and repair = 6.

Primary paired outcomes and contingency counts: For the 266 complete pairs the validator outcomes produced the following contingency counts: n11 = 260 (both baseline and guarded pass), n10 = 5 (guarded pass only), n01 = 0 (baseline pass only), and n00 = 1 (both fail). These counts yield observed per‑condition proportions baseline = 260/266 = 0.9774436090 and guarded = 265/266 = 0.9962406015; the paired (guarded - baseline) difference = 0.01879699248. The exact McNemar two‑sided test applied to the discordant pair counts (n10 = 5, n01 = 0) produced p = 0.0625. Because the exact two‑sided McNemar p for observed discordants is 0.0625 (1/16), it narrowly misses the preregistered \ensuremath{\alpha} = 0.05 rejection threshold.

Interpreting the small‑discordant pattern and the two inferential summaries: The discordant total (n10 + n01 = 5) is small. The exact two‑sided McNemar test calculates two‑sided probability mass under a Binomial(total=5, p=0.5) for outcomes as or more extreme than the observed (k >= 5 or k <= 0), yielding p = 2 * (0.5\textasciicircum{}5) = 0.0625. In contrast, the paired bootstrap percentile CI (B = 10,000 resamples; seed = 1729) for the paired difference is [0.0037593985, 0.0375939850] and excludes zero. Both summaries are valid descriptions of uncertainty; the discrete exact test is conservative when discordant counts are few, while the bootstrap CI summarizes empirical resampling variability. We therefore report both: the preregistered hypothesis test (McNemar) did not reject at \ensuremath{\alpha} = 0.05, while the paired bootstrap CI suggests a small positive effect whose lower bound is barely above zero for the observed sample.

Repair invocation, yield, and derived operational quantities: Repair was invoked in 6 episodes (stage\_counts: repair = 6). The contingency pattern (n10 = 5) indicates that at most five of those repair invocations contributed to additional guarded successes among the 266 analysed pairs; the run-level accounting (6 repair calls vs 5 discordant guarded wins) is consistent with a small number of successful repairs and at most one repair that did not produce a guarded pass among the complete pairs. Median repair iterations per task across complete pairs is 0 (most tasks did not require repair). A useful derived metric is the number needed to treat (NNT) interpretation: 1 / absolute\_risk\_reduction = 1 / 0.01879699248 \ensuremath{\approx} 53.2, i.e., roughly 53 guarded tasks (under this run's conditions and shared\_initial\_candidate semantics) were associated with one additional verifier‑approved configuration relative to baseline. This NNT is a descriptive operational scalar and should be interpreted in context: with a very high baseline pass rate, the marginal correction yield per repair opportunity is small.

Representative repair case diagnostics: Representative artifacts illustrate typical failure and repair behavior. Task-000008 (difficulty: "extreme", required\_command\_count = 9) provides a concrete example: the shared initial candidate contained a truncated final token line ("interface uplink2\textbackslash{}n") and failed the validator's parsed\_command\_count vs required\_command\_count check; the automated repair call produced a corrected candidate that completed the previously truncated command and included the required edge1 sequence. The guarded final validator trace for task-000008 shows validation\_after.valid = true and violation\_count = 0. These per‑episode validator traces record parsed\_command\_count, required\_command\_count, observed\_touched\_field\_count, normalized\_configuration, and violation codes, enabling targeted audit of repair logic and verifier coverage for each repaired episode.

Sensitivity to missingness and preregistered conservative handling: Thirty‑four both‑arm episodes were excluded from the primary confirmatory analysis (complete\_pairs = 266) per preregistered rules. A conservative sensitivity analysis that treats those 34 excluded both‑arm episodes as failures (complete N = 300) yields baseline = 260/300 = 0.866667 and guarded = 265/300 = 0.883333; the discordant counts among the complete pairs remain n10 = 5, n01 = 0, producing the same exact McNemar p = 0.0625 for the paired test on the complete‑pair subset. This shows the primary inferential conclusion (nonrejection under exact McNemar) is robust to preregistered conservative missingness treatment, while the absolute per‑condition proportions change due to inclusion of excluded episodes as failures.

Quantitative artifact summaries and reproducibility metadata used in Results: All numerical summaries reported above derive directly from archived analysis artifacts: the paired contingency counts, condition summaries, missingness summary, and provider‑call accounting produced the reported totals (complete\_pairs = 266; baseline\_success\_count = 260; guarded\_success\_count = 265; hosted\_model\_call\_event\_count = 306; repair calls = 6; both\_missing\_pairs = 34). The paired bootstrap used B = 10,000 resamples with seed = 1729 (bootstrap parameters archived) to produce the reported percentile CI. Per‑episode validator traces and scoring JSONs include normalized\_configuration, parsed\_command\_count, required\_command\_count, and violation codes for each episode and can be used to reproduce failure‑mode breakdowns or to recalibrate validator rule coverage.

Limitations specific to the observed Results: (1) High baseline pass rate produced few discordant pairs: rare failures reduce the statistical power to detect small absolute improvements under the preregistered exact McNemar test. (2) Repair calls were infrequent (6/300) in this corpus; larger absolute improvements from GVR are more plausible on corpora with lower first‑pass accuracy. (3) The NNT and repair yield reported are conditional on shared\_initial\_candidate semantics (the same initial sample across arms) and the particular model sampling/settings used; different sampling regimes or constrained first‑pass prompts could change the yield. (4) Per‑episode latency was not archived and thus the preregistered latency estimand could not be evaluated; timing instrumentation is recommended for future runs. (5) Validator coverage is finite: violations outside deterministic\_netops\_validator\_v5's rule battery are unmeasured here.

\section{Discussion}
The experiment produced a small absolute improvement in verifier pass‑rate when the allowed repair stage could operate on a shared initial candidate: guarded - baseline = 0.01879699248 (n11=260, n10=5, n01=0, n00=1). Practically, this pattern---few discordant pairs with all discordance in the guarded\ensuremath{\rightarrow}pass direction---indicates that (a) the initial LLM candidate was correct for most tasks in this corpus, and (b) the bounded repair stage occasionally corrected defects that the initial candidate incurred. The artifact accounting (repair calls = 6; hosted\_model\_call\_event\_count = 306) confirms repairs were rare relative to total generation activity, explaining the modest absolute gain.

From a statistical perspective the run highlights an important small‑sample inference nuance that directly affects operational interpretation. Exact McNemar (two‑sided) returned p = 0.0625 under the preregistered \ensuremath{\alpha}=0.05 decision rule, while a paired bootstrap percentile CI (B=10,000, seed=1729) produced a 95\% interval that narrowly excludes zero. These outcomes are consistent: with very few discordant pairs (here 5 vs 0), the exact two‑sided McNemar test is discrete and can be conservative; bootstrap resampling of paired differences summarizes empirical variability and can produce a CI excluding zero even when the exact discrete p slightly exceeds 0.05. For readers and practitioners this means that reporting both formal hypothesis outcomes and interval estimates (as archived here) gives a clearer picture: the effect magnitude is small but the CI suggests the effect is positive with narrow uncertainty in the observed corpus, whereas the exact hypothesis test under a strict \ensuremath{\alpha} threshold does not claim formal rejection.

Operationally, artifact evidence supports three practical inferences grounded in this run. First, when first‑pass generator accuracy is high (\ensuremath{\approx}97.7\% baseline here), bounded automated repair offers limited marginal benefit because few failures remain to correct; this is exactly what we observed (6 repair calls on 300 tasks). Second, the marginal utility of GVR will be larger when (i) first‑pass accuracy is lower, (ii) verifier coverage detects failure modes common in the target workload, or (iii) the repair mechanism is allowed more than one iteration---conditions not present in this bounded run. Third, the design of the verifier strongly conditions measured benefit: deterministic\_netops\_validator\_v5 enforces a finite, archived battery of syntactic/sequence/constraint/reachability checks; any semantic failure mode outside that battery is neither detected nor credited to the GVR pipeline. In short, GVR's measured improvement depends both on how often the generator errs and on whether the verifier can detect the errors that matter operationally.

Threats to inference and external validity remain artifact‑grounded. The preregistered shared\_initial\_candidate estimand isolates the effect of validator‑triggered repair for a fixed initial sample; it does not measure improvements obtainable by engineering the first pass (e.g., constrained prompting, constrained decoding, or multiple independent draws). Our execution manifest and reconciliation artifacts explicitly record stage\_counts (shared\_initial\_generation = 300, repair = 6) and the lack of independent first‑pass generations for separate arms, so extrapolating to comparisons that change initial sampling is unsupported by these artifacts. Missingness also reduced effective sample size: 34 both‑arm provider failures were excluded per preregistered rules, lowering precision and power; deterministic\_reconciliation.json and analysis/missingness\_summary.csv archive these counts. Finally, per‑episode latency was not recorded in per‑task artifacts and therefore the preregistered latency secondary estimand could not be evaluated---this absence constrains operational claims about responsiveness or real‑time applicability.

Design and research recommendations informed by these artifacts. (1) For stronger causal evidence about prompting or sampling strategies, future preregistrations should include additional independent arms that draw separate initial candidates (independent\_generation semantics) or include constrained‑first‑pass prompting; these require separate initial model calls per arm and differ from the shared\_initial\_candidate estimand used here. (2) To assess operational trade‑offs, pipeline runs must archive per‑episode timing (generation\ensuremath{\rightarrow}verification\ensuremath{\rightarrow}repair\ensuremath{\rightarrow}final validation) so latency and throughput can be measured; our provider accounting shows model call counts and failures but not per‑episode latencies. (3) Because verifier coverage materially shapes measured benefit, future studies should expand the verifier rule set and publish clear catalogs of detectable failure codes; the per‑episode validator traces archived in execution/scoring/*.json enable such audits and should be exploited to extend verifier coverage and to quantify verifier false negatives. (4) When baseline pass rates are high, consider targeting task corpora with intentionally lower baseline accuracy (e.g., out‑of‑distribution intents, noisier prompts, or reduced grounding) to increase discordant counts and thereby statistical power to detect repair benefits.

Concluding synthesis: the archived artifacts show that a bounded GVR pipeline can correct occasional deterministic‑validator failures for a fixed initial candidate, but when the generator's first‑pass accuracy is already high and the verifier's coverage is finite, the absolute operational improvement is small. The run's reconciled accounting, validator traces, repair provider traces, and contamination checks (no exact public duplicates flagged) provide a reproducible empirical substrate for these conclusions and for follow‑on experiments that expand sampling semantics, verifier coverage, or repair budgets.

\section{Limitations}
\begin{itemize}
\item Shared\_initial\_candidate semantics: both arms used the same single sampled initial candidate per preregistration; the estimand measures validator‑triggered repair benefit for that fixed initial sample and does not evaluate independent first‑pass generations or constrained first‑pass prompting strategies.
\item Missingness and sample reduction: planned N=300 pairs; after infrastructure/provider exclusions 266 pairs were complete and 34 both‑arm episodes were excluded per preregistered rules, reducing effective sample size and precision.
\item Small discordant counts: primary inference used exact McNemar with few discordants (n10=5, n01=0); conclusions are sensitive to inferential choice and limited power.
\item Validator coverage: deterministic\_netops\_validator\_v5 enforces a finite set of syntactic/semantic/reachability rules; semantic violations outside its coverage are possible and unmeasured.
\item Runtime nondeterminism and sampling: archived execution/model\_configuration.json records temperature = null and no deterministic seed; nondeterministic sampling cannot be excluded for recorded model calls.
\item H2 latency not evaluated: per‑episode end‑to‑end latency was not present in archived per‑task artifacts and therefore the preregistered latency bound (median <= 60 s) could not be assessed.
\item External validity: task corpus derives from public NetConfEval‑style examples and a seeded synthetic generator limited to two vendor‑dialect clusters; results do not imply universal benefit across other vendors, larger topologies, or production deployments.
\end{itemize}

\section{Conclusion}
We executed a preregistered paired experiment that evaluated a GVR pipeline (gpt‑5‑mini + deterministic\_netops\_validator\_v5 + <=1 automated repair) on intent\ensuremath{\rightarrow}CLI tasks under shared\_initial\_candidate semantics. Across 266 analysed pairs the guarded pipeline improved validator pass rate by 0.01879699248 but did not meet the preregistered exact‑McNemar \ensuremath{\alpha}=0.05 decision rule (p=0.0625); a paired bootstrap CI (B=10,000, seed=1729) narrowly excludes zero. Repair calls were rare (6/300) and median repair iterations per task = 0. Per‑episode latency was not archived and could not be assessed. All execution and analysis artifacts (raw responses, validator traces, provider call accounting, and reconciliation manifests) are archived in the evidence bundle to support reanalysis and follow‑on work.

References
[1] C. Wang, M. Scazzariello, A. Farshin, S. Ferlin, D. Kostić, and M. Chiesa, "NetConfEval: Can LLMs Facilitate Network Configuration?," in Proc. ACM, 2024, doi: 10.1145/3656296.
[2] S. Y. Long, J. Tan, B. Mao, F. Tang, Y. Li, M. Zhao, and N. Kato, "A Survey on Intelligent Network Operations and Performance Optimization Based on Large Language Models," IEEE Commun. Surveys \& Tutorials, 2025, doi: 10.1109/comst.2025.3526606.
[3] R. Cadenas, "Convergent Correctness in Stochastic Code Generation: A Generate‑Verify‑Repair Architecture for Deterministic Validation of Autonomous AI Coding Agents in Regulated Environments," SSRN, 2026, doi: 10.2139/ssrn.6754899.
[4] M. Brown, A. Fogel, D. Halperin, V. Heorhiadi, R. Mahajan, and T. Millstein, "Lessons from the evolution of the Batfish configuration analysis tool," Proc., 2023, doi: 10.1145/3603269.3604866.
[5] R. Mondal, A. Tang, R. Beckett, T. Millstein, and G. Varghese, "What do LLMs need to Synthesize Correct Router Configurations?," Proc., 2023, doi: 10.1145/3626111.3628194.

\section*{Disclosure Statement}
Disclosure (concise): Declared model: gpt-5-mini (execution/model\_configuration.json archived; declared\_model\_version = "gpt-5-mini"). Orchestration adapter: hosted\_netops\_gvr\_v1. Deterministic validator: deterministic\_netops\_validator\_v5. Preregistration: cnsms2026-001-gvr-batfish-prereg-v1 (manifest SHA‑256: 46a22349\allowbreak{}2f760950\allowbreak{}b3d06f47\allowbreak{}5129ce21\allowbreak{}e52e296c\allowbreak{}4c2a1e09\allowbreak{}df3ccd50\allowbreak{}bec9f891). Initial master prompt immutable reference: provenance/master\_prompt.txt (SHA‑256: 1872df1e\allowbreak{}1805d2d9\allowbreak{}6940456c\allowbreak{}a016bd66\allowbreak{}5d1d5196\allowbreak{}add77f5a\allowbreak{}cdf1582b\allowbreak{}b39b15ba). Human role and autonomy boundary: before the autonomy lock a human Research Director selected the topic family and approved pre‑lock infrastructure development; after the autonomy lock the registered pipeline autonomously performed literature selection, preregistration, experimental execution, analysis, manuscript drafting, AI peer review simulation, and revision. Nondeterminism disclosure: archived execution/model\_configuration.json records temperature = null (unset) and no deterministic sampling seed is archived; therefore nondeterministic sampling of model outputs cannot be excluded for the recorded provider calls. Execution and analysis artifacts, logs, and hash manifests are archived in the evidence bundle referenced in the preregistration.

\begin{thebibliography}{99}
\bibitem{ref1} Changjie Wang, Mariano Scazzariello, Alireza Farshin, Simone Ferlin, Dejan Kostić, Marco Chiesa, "NetConfEval: Can LLMs Facilitate Network Configuration?", 2024, 10.1145/3656296
\bibitem{ref2} S.Y. Long, Jingjing Tan, Bomin Mao, Fengxiao Tang, Yangfan Li, Ming Zhao, Nei Kato, "A Survey on Intelligent Network Operations and Performance Optimization Based on Large Language Models", 2025, 10.1109/comst.2025.3526606
\bibitem{ref3} Rafael Cadenas, "Convergent Correctness in Stochastic Code Generation: A Generate-Verify-Repair Architecture for Deterministic Validation of Autonomous AI Coding Agents in Regulated Environments", 2026, 10.2139/ssrn.6754899
\bibitem{ref4} Matt Brown, Ari Fogel, Daniel Halperin, Victor Heorhiadi, Ratul Mahajan, Todd Millstein, "Lessons from the evolution of the Batfish configuration analysis tool", 2023, 10.1145/3603269.3604866
\bibitem{ref5} Rajdeep Mondal, Alan Tang, Ryan Beckett, Todd Millstein, George Varghese, "What do LLMs need to Synthesize Correct Router Configurations?", 2023, 10.1145/3626111.3628194
\end{thebibliography}

\end{document}
